\documentclass[10pt,a4paper]{amsart}
\usepackage[margin=19mm]{geometry}
\usepackage{amsmath,amssymb,mathtools}
\usepackage[scr=rsfs,cal=euler]{mathalfa}
\usepackage{xfrac}
\usepackage[unicode,hidelinks,bookmarks=false]{hyperref}
\newcommand{\ie}{i.e.\ }
\newcommand{\cf}{cf.\ }
\newcommand{\resp}{respectively\ }
\newcommand{\Subsection}{\subsection}
\providecommand{\nobreakdash}{\nobreak\hbox{-}}
\setlength{\emergencystretch}{2em}
\multlinegap0pt
\usepackage{amssymb}
\newtheorem{lemma}{Lemma}
\newcommand{\C}{\mathbb C}
\title{A two-dimensional counterexample for Kobayashi completeness of balanced domains}
\author{Armen Edigarian}
\date{Source: arXiv:2609.07697v1 (CC BY 4.0)}
\begin{document}
\maketitle

\noindent\emph{Source and licence.} A two-dimensional counterexample for Kobayashi completeness of balanced domains. Version arXiv:2609.07697v1 (CC BY 4.0), distributed by arXiv under the Creative Commons Attribution 4.0 International licence, http://creativecommons.org/licenses/by/4.0/, as recorded in the arXiv licence metadata for this exact version and shown on the abstract page https://arxiv.org/abs/2609.07697v1. Full licence notice: \texttt{licenses/arxiv-2609.07697-CC-BY-4.0.txt}.\par\smallskip

\noindent\textit{Editorial scope.} Counted unit: the article's lemma lem:sublevel (source0.tex lines 96--104). The article's complete original proof of it is delivered with it (source0.tex lines 106--135, one \texttt{proof} environment of 30 lines). No mathematics is written, repaired or re-derived: the statement and the proof are the article's own lines, in the article's own notation, and no retained line is substituted or re-flowed.  No further source-local context is needed: every cross-reference of every retained line resolves inside the two delivered spans.  Nothing was omitted from any retained span, so the package carries no omission record.  No retained line contains a citation, so the wrapper carries no bibliography and no bibliography entry.  No retained line contains a figure environment, an image-inclusion command or drawing code, so no image is delivered and the package declares no asset dependency.  Editorial formatting only, in addition: an AMS \texttt{amsart} preamble that restates the article's own declarations and macros used by the retained lines, namely the environments \texttt{lemma} on separate counters, together with the standard packages those lines need.  The retained environments print in this document as Lemma 1 in order of appearance, whereas the article prints the same items with the numbering of its own full document.  The complete native source of the article as distributed in the arXiv e-print for this exact version is bundled unchanged as \texttt{sources/209651/source0.tex}.
\begin{lemma}\label{lem:sublevel}
For every $\zeta\in\mathbb C$ and every $j\ge1$,
\[
 \psi_j(\zeta)\le
 \frac{m_j}{a_j(1+\tau_j^2)}\,|\zeta|.
\]
Consequently,
$u(\zeta)\le \frac45|\zeta|$ for any $\zeta\in\mathbb C$.
\end{lemma}

\begin{proof} Note that
$\log(1+x+\frac{x^2}{2})\le x\quad\text{ for any }x\ge0$. Since $0<\tau_j<1$, we have
\[
 \psi_j(\zeta)
 \le \frac{m_j}{2}
 \log\frac{(1+\frac{|\zeta|}{a_j})^2+\tau_j^2}{1+\tau_j^2}\le \frac{m_j}{a_j(1+\tau_j^2)}\,|\zeta|.
\]
We now estimate $u$. We have  
$\log(1+|\zeta|^2)\le |\zeta|$ for any $\zeta\in\C$.
Using the estimate for each $\psi_j$ we obtain
$$ 
u(\zeta)
 \le \left(
 \frac{7}{15}
 +\sum_{j=1}^{\infty}
 \frac{m_j}{a_j(1+\tau_j^2)}
 \right)|\zeta|\le \left(
 \frac{7}{15}
 +\sum_{j=1}^{\infty}\frac{m_j}{a_j}
 \right)|\zeta|.
$$
We have $\sum_{j=1}^{\infty}\frac{m_j}{a_j}
 =\sum_{j=1}^{\infty}4^{-j}=\frac13$.
Consequently
\[
 u(\zeta)\le
 \left(\frac{7}{15}+\frac13\right)|\zeta|
 =\frac45|\zeta|.
\]
\end{proof}

\end{document}
